\documentclass[12pt,a4paper]{letter}

\usepackage[french]{babel}
\usepackage[T1]{fontenc}
\usepackage[utf8]{inputenc}
\usepackage{lmodern}
\usepackage{geometry}
\usepackage[hidelinks]{hyperref}

\geometry{a4paper, top=2.5cm, bottom=2.5cm, left=3cm, right=3cm}

\signature{Eriam Schaffter}
\address{%
  Eriam Schaffter\\
  Enseignant agrégé ISI --- BTS SIO\\
  Chargé de cours M2 MIAGE, Université Claude Bernard Lyon~1\\
  \href{mailto:eriam@mediavirtuel.com}{eriam@mediavirtuel.com}
}

\begin{document}

% --- VERSION POUR SIM ---
% Décommenter et adapter selon la revue cible

\begin{letter}{%
  Comité éditorial\\
  \textit{Systèmes d'Information et Management}\\
  % ou : \textit{Revue Française de Gestion}
}

\opening{Madame, Monsieur,}

J'ai l'honneur de soumettre à votre revue un article intitulé \textbf{\enquote{L'expert-comptable comme tiers garant du code opposable : analyse prospective d'une recomposition professionnelle à l'ère des mutations logicielles et réglementaires}}.

% --- POUR SIM ---
Cet article s'inscrit pleinement dans le périmètre éditorial de \textit{Systèmes d'Information et Management}. Il aborde la question de la gouvernance des systèmes d'information des PME dans un contexte de mutation technologique (émergence des grands modèles de langage) et de durcissement réglementaire (RGPD, AI Act, NIS~2, directive 2024/2853). Il développe l'hypothèse selon laquelle la profession d'expert-comptable, par ses actifs institutionnels spécifiques, est structurellement positionnée pour devenir le tiers garant de la conformité numérique des PME françaises.

% --- POUR RFG (remplacer le paragraphe ci-dessus) ---
% Cet article s'inscrit dans le périmètre éditorial de la \textit{Revue Française de Gestion}
% sur l'axe de la transformation des professions réglementées face aux mutations
% technologiques. Il analyse, en mobilisant la sociologie des professions (Abbott, 1988)
% et la théorie économique des tiers de confiance, les conditions dans lesquelles
% la profession d'expert-comptable peut repositionner son périmètre d'activité vers
% la certification de la conformité des systèmes d'information des PME.

Sur le plan académique, l'article contribue à l'articulation de thématiques rarement traitées conjointement dans la littérature : la sociologie des professions réglementées, la gouvernance des SI dans les PME, la dimension normative du code informatique et la théorie économique des tiers de confiance. Sur le plan pratique, il intéresse directement les décideurs de la profession et les organismes de formation face à une transformation de grande ampleur.

L'article est présenté comme un essai théorique argumenté à vocation prospective, conformément aux normes éditoriales pour ce type de contribution en sciences de gestion. Il comporte environ [XX] mots (hors bibliographie et résumés), un résumé en français et en anglais avec mots-clés, et [XX] références bibliographiques.

Je certifie que cet article est original et n'est pas soumis simultanément à une autre revue. Je n'ai pas de conflit d'intérêts à déclarer en rapport avec ce travail.

Je me tiens à votre disposition pour tout renseignement complémentaire et reste dans l'attente de votre réponse.

\closing{Dans l'attente de votre retour, je vous adresse mes sincères salutations,}

\end{letter}

\end{document}
